\documentclass[12pt,a4paper]{article}
\usepackage[T2A]{fontenc}
\usepackage[utf8]{inputenc}
\usepackage[russian,english]{babel}
\usepackage{amsmath}
\usepackage{booktabs}
\usepackage{geometry}
\usepackage{hyperref}
\geometry{margin=2.5cm}
\hypersetup{hidelinks}

\title{StableToolBenchQoS: моделирование качества работы API}
\author{}
\date{}

\begin{document}
\maketitle

\section{Цель проекта}

Функционально похожие инструменты могут различаться вероятностью успешного
выполнения, временем ответа и ценой. Если агенту доступно несколько подходящих
API, выбор только по семантическому соответствию не позволяет учитывать эти
различия.

Цель проекта StableToolBenchQoS --- создать воспроизводимую среду, в которой
инструменты ToolBench имеют три дополнительные характеристики:
\begin{itemize}
    \item \texttt{success\_rate} --- вероятность успешного вызова;
    \item \texttt{latency} --- время выполнения вызова;
    \item \texttt{cost} --- синтетическая денежная стоимость вызова.
\end{itemize}

Эта среда нужна для оценки алгоритма выбора среди функционально близких
инструментов. В частности, она позволяет проверить, способен ли ToolBandit
повысить долю успешно решённых заданий и одновременно соблюдать ограничения
на бюджет и время ответа.

\section{Граф функциональных отношений между API}

Чтобы сравнивать стратегии выбора инструмента, сначала необходимо определить,
какие API могут решать одну и ту же задачу. Для этого в проекте строится граф
функциональных отношений. Его исходные вершины --- $6\,470$ API из рабочего
каталога. Пара API может получить одно из трёх отношений:
\begin{itemize}
    \item \texttt{interchangeable} --- API решают одну пользовательскую задачу
    и могут заменять друг друга; необходимые преобразования аргументов не
    требуют обращения к внешним данным;
    \item \texttt{contains} --- API имеют общую основную функцию, но один из
    них предоставляет более широкие возможности, область данных или набор
    результатов;
    \item \texttt{different\_capability} --- взаимозаменяемость или включение
    по документации не подтверждаются.
\end{itemize}

Полный перебор всех пар потребовал бы проверить более $20$ миллионов
комбинаций. Поэтому сначала выполняется поиск кандидатов. Текстовые описания,
названия и параметры API кодируются моделью
\texttt{Qwen/Qwen3-Embedding-0.6B}; дополнительно используется лексический
поиск BM25. Для каждого API сохраняются до 30 ближайших соседей каждого метода,
после чего совпадающие неупорядоченные пары объединяются.

Qwen выбран как практический компромисс для этого этапа. Модель строит
instruction-aware embeddings: при кодировании можно явно указать, что требуется
найти API с той же или близкой функцией, а не просто текст с похожими словами.
Контекст до 2048 токенов позволяет учитывать не только название endpoint, но и
описание, параметры и свойства вызова. Размер модели $0{,}6$ млрд параметров
позволяет локально обработать весь каталог на одной видеокарте, не отправляя
документацию внешнему сервису и не оплачивая каждый embedding. Это был
инженерный выбор по сочетанию семантического поиска, длины контекста и стоимости
вычислений; отдельного сравнения нескольких embedding-моделей в проекте пока
не проводилось.

BM25 добавлен потому, что embedding-поиск может пропускать пары с редкими
названиями методов, одинаковыми названиями параметров или точными техническими
терминами. Таким образом, Qwen отвечает прежде всего за семантическую близость,
а BM25 дополняет её точным лексическим совпадением.

В результате dense-поиск нашёл $132\,182$ пары, BM25 --- $141\,439$ пар;
$54\,806$ пар были найдены обоими методами. После объединения сформирован пул
из $218\,815$ уникальных пар-кандидатов. Embedding score применяется для
ранжирования кандидатов, но сам по себе не считается меткой функционального
отношения.

На следующем этапе LLM получает документацию двух API, включая их описания и
параметры, и выбирает одно из трёх отношений. Для \texttt{contains} модель
также указывает направление, то есть какой API функционально шире. Вместе с
решением сохраняются confidence, краткое обоснование, свидетельства из
документации и признак необходимости ручной проверки. Основной запуск выполнил
$14\,000$ LLM-вызовов. К нему были добавлены $287$ решений из предыдущих
экспериментов, а ещё $2\,647$ решений выведены
из уже известных отношений без нового обращения к модели. Журнал до ручных
исправлений содержит $16\,934$ записи, из которых $14\,287$ являются прямыми
решениями LLM или перенесёнными ранее решениями.

Кластеры строятся только по отношению \texttt{interchangeable}: все связанные
им API объединяются в одну компоненту эквивалентности. Отношение
\texttt{contains} сохраняется как направленное ребро между получившимися
кластерами, а \texttt{different\_capability} --- как свидетельство различия
между ними. После применения ручных исправлений итоговый граф содержит
$6\,027$ кластеров и $13\,773$ межкластерных ребра. Из кластеров $5\,753$
содержат один API, $274$ --- несколько API, а максимальный размер кластера
равен 19.

Граф является частичным: отсутствие ребра не означает, что два API точно не
связаны. Оно может означать, что пара не попала в пул кандидатов или ещё не
была размечена. Для экспериментального набора используются только явно
построенные кластеры взаимозаменяемых API.

\section{Моделирование QoS}

\subsection{Расширение виртуального сервера}

В форке сохранена схема виртуального сервера StableToolBench: при успешном
вызове ответ берётся из официального или дополнительного кеша, а при промахе
генерируется LLM. Перед получением ответа сервер выполняет моделирование QoS.
С вероятностью $1-p_a$, где $p_a$ --- \texttt{success\_rate} API $a$, сервер
возвращает ошибку вызова и не обращается к LLM. Это моделирует недоступность API
и предотвращает лишние платные запросы к симулятору.

Задержка отдельного вызова моделируется логнормальной случайной величиной:
\[
L_a \sim \operatorname{LogNormal}(\mu_a,\sigma_a^2).
\]
Параметр $\mu_a$ выбирается так, чтобы математическое ожидание распределения
совпадало с заданной для API величиной \texttt{expected\_latency\_ms}.
Логнормальное распределение гарантирует положительность задержки и допускает
редкие медленные вызовы. Сервер может как реально ожидать сгенерированное
время, так и только возвращать его в метаданных.

Ответ виртуального сервера дополнен объектом \texttt{qos}, содержащим сценарий,
результат вызова, вероятность успеха, сгенерированную и ожидаемую задержки,
стоимость и порядковый номер вызова.

\subsection{Формирование QoS-профилей}

Актуальный набор содержит $7\,546$ QoS-профилей. Он покрывает все $2\,491$ API
из solvable-набора и все $6\,470$ API из графа отношений; $1\,415$ API входят
в оба множества. Для каждого API создаётся базовый профиль.

Распределения \texttt{success\_rate} и ожидаемой задержки строятся методом
воспроизводимой bootstrap-выборки из исходных метрик StableToolBench. API
получает метрики другого инструмента-донора, выбранного по идентификатору API и
\texttt{seed}. Базовая вероятность успеха вычисляется как
\[
p_a = \frac{\mathrm{avgServiceLevel}}{100}
      \cdot
      \frac{\mathrm{avgSuccessRate}}{100}.
\]
Ожидаемая задержка берётся из поля \texttt{avgLatency} того же донора. Общий
донор сохраняет наблюдаемую связь между доступностью и задержкой, но не
приписывает endpoint его собственную историческую метрику.

Числовой цены в исходном StableToolBench нет, поэтому стоимость генерируется
синтетически из логнормального распределения. В текущей конфигурации она
ограничена диапазоном от $0{,}00005$ до $0{,}05$ условных денежных единиц за
вызов.

Из базового профиля формируются три сценария:
\begin{table}[ht]
    \centering
    \begin{tabular}{lccc}
        \toprule
        Сценарий & Success rate & Ожидаемая задержка & Стоимость \\
        \midrule
        normal   & $p_a$         & $L_a$  & $c_a$ \\
        degraded & $0{,}6375p_a$ & $2L_a$ & $c_a$ \\
        outage   & $0{,}01p_a$   & $4L_a$ & $c_a$ \\
        \bottomrule
    \end{tabular}
    \caption{Сценарии моделирования QoS.}
\end{table}

Сценарий задаётся один раз при запуске сервера и применяется ко всем вызовам в
эксперименте. При одинаковых конфигурации, \texttt{seed} и порядке запросов
последовательность успехов, отказов и задержек воспроизводится.

\section{Кеш ответов и набор заданий}

Дополнительный кеш строится для API, используемых в заданиях. Для API с
параметрами целевое количество равно трём различным примерам, а для API без
параметров достаточно одного. Для актуального кластерного набора требуется
$918$ API: $914$ из них достигают целевого покрытия, а для четырёх полный набор
примеров пока не собран. Это соответствует покрытию около $99{,}6\%$ и не
блокирует работу, поскольку при промахе кеша сервер может обратиться к
LLM-симулятору.

Отдельно построен граф функциональных отношений между $6\,470$ API. После
применения ручных исправлений он содержит $6\,027$ кластеров, из которых $274$
содержат более одного API. Максимальный размер кластера равен $19$.

Для подготовки эксперимента существующие естественные запросы из
solvable-набора сопоставляются с графом. Если исходный API входит в кластер с
несколькими элементами, в список доступных инструментов добавляются остальные
API этого кластера. Из $765$ исходных тестовых заданий этому условию
соответствуют $193$. Они охватывают $119$ кластеров; суммарно в списки
кандидатов добавлен $991$ альтернативный API.

Для каждого из $193$ заданий с помощью LLM создано по $20$ перефразировок.
Вместе с исходными формулировками расширенный набор содержит $4\,053$ задания.
При разбиении этого набора на обучающую и тестовую части все варианты одного
исходного задания следует помещать в одну часть, иначе почти одинаковые
формулировки могут попасть одновременно в обучение и тестирование.

\subsection{Ручная проверка графа}

Для проверки отношений между API был подготовлен интерфейс ручной разметки.
В журналах разметки содержится $302$ уникальные проверенные пары. С текущими
решениями графа удалось сопоставить $242$ пары: для $229$ из них решение модели
было подтверждено, а для $13$ изменено. Ещё $60$ записей не были сопоставлены с
текущим набором решений и поэтому не применялись при построении графа.

После применения ручных исправлений граф содержит $13\,773$ ребра: $825$ с
отношением \texttt{contains} и $12\,948$ с отношением
\texttt{different\_capability}. В метаданных также остаются $16$
диагностических записей, которые следует учитывать при дальнейшей ревизии
графа.

\subsection{Результаты анализа разметки}

Отдельный анализ выполнен для $14\,287$ прямых решений исходного графа до
применения ручных исправлений. В это число входят $14\,000$ решений основного
запуска и $287$ решений, перенесённых из предыдущих экспериментов. Для всех
этих пар доступен embedding score.

\begin{table}[ht]
    \centering
    \small
    \begin{tabular}{lrrrr}
        \toprule
        Embedding score & Всего & Interchangeable & Contains & Different \\
        \midrule
        $[0{,}85; 0{,}90)$ & $6\,866$ & 93  & 146 & $6\,627$ \\
        $[0{,}90; 0{,}95)$ & $4\,960$ & 187 & 292 & $4\,481$ \\
        $[0{,}95; 1{,}00]$ & $2\,461$ & 175 & 385 & $1\,901$ \\
        \midrule
        Всего                & $14\,287$ & 455 & 823 & $13\,009$ \\
        \bottomrule
    \end{tabular}
    \caption{Решения модели в зависимости от embedding score.}
\end{table}

BM25 score доступен для $12\,956$ из $14\,287$ прямых решений. Поскольку
значения BM25 зависят от корпуса и длины текста и не нормированы на интервале
от нуля до единицы, для анализа использованы квартильные диапазоны. Если пара
была найдена в обоих направлениях, её score определяется как максимальное из
двух значений.

\begin{table}[ht]
    \centering
    \small
    \begin{tabular}{lrrrr}
        \toprule
        BM25 score & Всего & Interchangeable & Contains & Different \\
        \midrule
        $[8{,}39; 52{,}09)$    & $3\,239$ & 225 & 173 & $2\,841$ \\
        $[52{,}09; 91{,}57)$   & $3\,190$ & 106 & 164 & $2\,920$ \\
        $[91{,}57; 144{,}91)$  & $3\,287$ & 46  & 280 & $2\,961$ \\
        $[144{,}91; 491{,}70]$ & $3\,240$ & 50  & 184 & $3\,006$ \\
        \midrule
        Всего                   & $12\,956$ & 427 & 801 & $11\,728$ \\
        \bottomrule
    \end{tabular}
    \caption{Решения модели по квартилям BM25 score.}
    \label{tab:decisions-by-bm25}
\end{table}

Для оценки ошибок использовались $222$ уникальные пары из очереди проверки
графа, для которых имеется ручная разметка. В $12$ случаях решение модели не
совпало с решением человека, что составляет $5{,}4\%$ проверенной выборки.
Распределение ручных решений и ошибок по embedding score приведено в
таблице~\ref{tab:review-by-embedding}.

\begin{table}[ht]
    \centering
    \small
    \begin{tabular}{lrrrrrr}
        \toprule
        Score & Проверено & Interch. & Contains & Different & Ошибок & Доля \\
        \midrule
        $[0{,}85; 0{,}90)$ & 40  & 36  & 4 & 0 & 4 & $10{,}0\%$ \\
        $[0{,}90; 0{,}95)$ & 61  & 58  & 1 & 2 & 3 & $4{,}9\%$ \\
        $[0{,}95; 1{,}00]$ & 121 & 116 & 1 & 4 & 5 & $4{,}1\%$ \\
        \midrule
        Всего                & 222 & 210 & 6 & 6 & 12 & $5{,}4\%$ \\
        \bottomrule
    \end{tabular}
    \caption{Ручная проверка по диапазонам embedding score.}
    \label{tab:review-by-embedding}
\end{table}

Из проверенных пар BM25 score доступен для $213$. Результаты по тем же
квартильным границам приведены в таблице~\ref{tab:review-by-bm25}.

\begin{table}[ht]
    \centering
    \small
    \begin{tabular}{lrrrrrr}
        \toprule
        BM25 score & Проверено & Interch. & Contains & Different & Ошибок & Доля \\
        \midrule
        $[8{,}39; 52{,}09)$    & 90 & 83 & 6 & 1 & 7 & $7{,}8\%$ \\
        $[52{,}09; 91{,}57)$   & 45 & 44 & 0 & 1 & 1 & $2{,}2\%$ \\
        $[91{,}57; 144{,}91)$  & 33 & 33 & 0 & 0 & 0 & $0{,}0\%$ \\
        $[144{,}91; 491{,}70]$ & 45 & 41 & 0 & 4 & 4 & $8{,}9\%$ \\
        \midrule
        Всего                   & 213 & 201 & 6 & 6 & 12 & $5{,}6\%$ \\
        \bottomrule
    \end{tabular}
    \caption{Ручная проверка по квартилям BM25 score.}
    \label{tab:review-by-bm25}
\end{table}

В проверенной выборке доля ошибок не меняется монотонно с ростом BM25 score:
она составляет $7{,}8\%$ в нижнем квартиле, $2{,}2\%$ и $0\%$ в двух средних
квартилях и $8{,}9\%$ в верхнем. Высокое лексическое сходство может возникать
как у действительно взаимозаменяемых API, так и у API с похожими названиями и
описаниями, но разными функциями. Поэтому BM25 полезен для отбора кандидатов,
но не заменяет функциональную классификацию пары.

В таблице~\ref{tab:review-by-confidence} те же проверки сгруппированы по
confidence, указанному LLM. Значение $0{,}78$ встречается только у двух пар,
поэтому долю ошибок $100\%$ для него нельзя интерпретировать как устойчивую
оценку качества.

\begin{table}[ht]
    \centering
    \begin{tabular}{rrrr}
        \toprule
        Confidence & Проверено & Ошибок & Доля ошибок \\
        \midrule
        $0{,}60$ & 2  & 0 & $0{,}0\%$ \\
        $0{,}70$ & 35 & 4 & $11{,}4\%$ \\
        $0{,}78$ & 2  & 2 & $100{,}0\%$ \\
        $0{,}80$ & 35 & 1 & $2{,}9\%$ \\
        $0{,}85$ & 38 & 1 & $2{,}6\%$ \\
        $0{,}90$ & 54 & 2 & $3{,}7\%$ \\
        $0{,}95$ & 55 & 2 & $3{,}6\%$ \\
        $1{,}00$ & 1  & 0 & $0{,}0\%$ \\
        \bottomrule
    \end{tabular}
    \caption{Расхождения с ручной разметкой по confidence LLM.}
    \label{tab:review-by-confidence}
\end{table}

Проверенная выборка не является случайной: при разметке можно было специально
отбирать пары с низкой уверенностью модели. Поэтому общую долю ошибок всех
$14\,287$ решений нельзя оценивать непосредственно по наблюдаемым $5{,}4\%$.
Полученные значения характеризуют только фактически проверенные пары.

\subsection{Ограничения}

Виртуальный ответ является моделью поведения API, а не результатом обращения к
реальному сервису. При точном попадании в кеш воспроизводится сохранённый
результат. При промахе качество зависит от LLM-симулятора: он может вернуть
правдоподобную структуру, но также способен сформировать слишком общий или
фактически неверный ответ. Поэтому эксперимент позволяет воспроизводимо
сравнивать стратегии маршрутизации, но не измеряет актуальное состояние
реальных сервисов.

Значения QoS также не являются текущими показателями конкретных API.
Вероятность успеха и ожидаемая задержка получены переносом исторического
профиля другого инструмента, а стоимость полностью синтетическая. Сценарии
\texttt{normal}, \texttt{degraded} и \texttt{outage} применяются ко всему
запуску, а не моделируют независимое изменение состояния каждого API во
времени.

Перефразировки одного задания не являются независимыми задачами. Их нельзя
случайно распределять между обучающей и тестовой выборками без группировки по
исходному заданию.

\begin{thebibliography}{9}

\bibitem{poon2025multillm}
Poon M., Dai X., Liu X., Kong F., Lui J.~C.~S., Zuo J.
Online Multi-LLM Selection via Contextual Bandits under Unstructured Context
Evolution. arXiv:2506.17670, 2025.
URL: \url{https://arxiv.org/abs/2506.17670}

\bibitem{chu2026lqm}
Chu K., Xiang D., Zhang W.
Latency-Quality Routing for Functionally Equivalent Tools in LLM Agents.
arXiv:2605.14241, 2026.
URL: \url{https://arxiv.org/abs/2605.14241}

\bibitem{qin2023toolllm}
Qin Y., Liang S., Ye Y. et al.
ToolLLM: Facilitating Large Language Models to Master 16000+ Real-world APIs.
ICLR 2024.
URL: \url{https://arxiv.org/abs/2307.16789}

\bibitem{toolbenchRepository}
OpenBMB. ToolBench: An Open Platform for Training, Serving, and Evaluating
Large Language Models for Tool Learning.
URL: \url{https://github.com/OpenBMB/ToolBench}

\bibitem{guo2024stabletoolbench}
Guo Z., Cheng S., Wang H. et al.
StableToolBench: Towards Stable Large-Scale Benchmarking on Tool Learning of
Large Language Models. Findings of ACL 2024.
URL: \url{https://aclanthology.org/2024.findings-acl.664/}

\end{thebibliography}
\end{document}
